\begin{abstract}
%随着大规模语言模型（LLMs）向实际部署迈进，微缩放格式 MXFP4（Microscaling FP4）凭借其兼顾高动态范围与硬件效率的特性，被视为下一代低比特推理的基石。然而，直接将 MXFP4 应用于 LLM 激活量化会导致显著的精度衰退。本文从理论层面剖析了 MXFP4 激活量化的误差结构，揭示了导致精度下降的根本原因在于激活分布与 MXFP4 块浮点结构之间存在的两种失衡：1.跨块方差极度不均衡：在按固定组大小进行 microscaling 时，组间方差高度不均衡, 少数高方差组主导总体量化误差并迫使共享尺度抬升，导致组内大量中小幅值激活被过度粗化。2.块内码本利用不均衡：真实激活的重尾分布与  MXFP4 几何级数码本严重失配，导致有效码字利用率极低，大量码字表示能力被闲置。
%为解决上述问题，我们提出了 TORQ (Two-level Orthogonal Rotation for MXFP4 Quantization) 。这是一种无需重训练的 PTQ 框架，旨在通过构建最优坐标变换来重塑激活空间的几何特性。 在宏观层面，TORQ 基于 Schur-Horn 定理，通过块间正交旋转实现方差的全局均衡，并从理论上证明了在凸误差假设下，该均衡配置等价于全局最小均方误差（MSE）。 在微观层面，TORQ引入最大熵原理，在组内执行正交旋转，最大化码字信息熵，使量化后激活的分布更适应fp数据分布，从而提升 MXFP4 码本的有效利用率。 在 LLaMA3和 Qwen3 等多种主流 LLM 上的实验表明，TORQ 相比现有 SOTA 方法显著提升了激活 MXFP4 PTQ 的精度：在Qwen3-32B模型上，在 WikiText 上将困惑度降至 8.43（BF16是7.61），在平均准确率上从 直接RTN的38.40% 提升至 73.63%(BF16是74.82)，成功弥合了 4-bit 浮点量化与全精度推理的鸿沟。

As Large Language Models (LLMs) advance toward practical deployment, the Microscaling FP4 (MXFP4) format has emerged as a cornerstone for next-generation low-bit inference, owing to its ability to balance high dynamic range with hardware efficiency. However, directly applying MXFP4 to LLM activation quantization inevitably leads to significant accuracy degradation. In this paper, we theoretically analyze the error structure of MXFP4 activation quantization, revealing that the root cause of this performance drop lies in two structural imbalances between activation distributions and the MXFP4 block floating-point format: (1).Extreme Inter-block Variance Imbalance: Under fixed-block microscaling, inter-block variance is highly uneven. Specifically, a minority of high-variance blocks dominate the total quantization error and force the shared scaling factor upward, resulting in the excessive coarsening of numerous small-magnitude activations within the block. 
(2).Intra-block Codebook Utilization Imbalance: The heavy-tailed distribution of real-world activations severely mismatches the MXFP4 geometric progression codebook, leading to extremely low effective codeword utilization and leaving significant representation capacity idle.
To address these challenges, we propose TORQ (Two-level Orthogonal Rotation for MXFP4 Quantization), a training-free Post-Training Quantization (PTQ) framework designed to reshape the geometric properties of the activation space through optimal coordinate transformations. 
% At the macroscopic level, TORQ leverages the Schur-Horn theorem to achieve global variance equalization via inter-block orthogonal rotation; we theoretically prove that this equilibrium configuration is equivalent to the global minimum Mean Squared Error (MSE) under convex error assumptions. 
At the macroscopic level, TORQ leverages the Schur-Horn theorem to redistribute activation energy via inter-block orthogonal rotation. Mathematically formalized as variance equalization, this process acts as a soft outlier suppression mechanism—preventing high-variance blocks from driving up the shared scaling factors and thereby preserving the precision of small-magnitude elements.
% At the microscopic level, TORQ introduces the principle of maximum entropy to execute intra-block orthogonal rotation. This maximizes codeword information entropy to better align the post-quantization activation distribution with the FP4 data distribution, thereby enhancing the effective utilization of the MXFP4 codebook. 
At the microscopic level, TORQ employs maximum-entropy-guided intra-block rotation to alleviate "codebook collapse." Actively dispersing the highly concentrated activation distribution, it effectively "activates" underutilized codewords to maximize the MXFP4 codebook’s information capacity.
Experiments on various mainstream LLMs such as LLaMA3 and Qwen3 show that TORQ significantly improves the accuracy of MXFP4 activation quantization compared to existing SOTA methods: on the Qwen3-32B model, the perplexity on WikiText is reduced to 8.43 (7.61 for BF16), and the average accuracy is increased from 38.40\% with direct RTN to 73.63\% (74.82\% for BF16), successfully bridging the gap between 4-bit floating-point quantization and full-precision inference.
\end{abstract}